\documentclass[10pt,a4paper]{amsart}
\usepackage[margin=19mm]{geometry}
\usepackage{amsmath,amssymb,mathtools}
\usepackage[scr=rsfs,cal=euler]{mathalfa}
\usepackage{xfrac}
\usepackage[unicode,hidelinks,bookmarks=false]{hyperref}
\newcommand{\ie}{i.e.\ }
\newcommand{\cf}{cf.\ }
\newcommand{\resp}{respectively\ }
\newcommand{\Subsection}{\subsection}
\providecommand{\nobreakdash}{\nobreak\hbox{-}}
\setlength{\emergencystretch}{2em}
\multlinegap0pt
\usepackage{bm}
\usepackage{bbm}
\usepackage{dsfont}
\usepackage{xspace}
\usepackage{cancel}
\usepackage{mathtools}
\usepackage[shortlabels]{enumitem}
\usepackage{amsthm}
\makeatletter
\@ifundefined{theorem}{\newtheorem{theorem}{Theorem}}{}
\@ifundefined{lemma}{\newtheorem{lemma}[theorem]{Lemma}}{}
\@ifundefined{prop}{\newtheorem{prop}[theorem]{Proposition}}{}
\makeatother
\newcommand{\delbar}{\overline{\partial}}
\newcommand{\del}{\partial}
\newcommand{\llangle}{\langle\!\langle}
\newcommand{\rrangle}{\rangle\!\rangle}
\title[The second variation of 2-spheres in the 2n-sphere]{The second variation of 2-spheres in the 2n-sphere}
\author{Gavin Ball, Jesse Madnick}
\date{Source: arXiv:2609.17887v1 (CC BY 4.0)}
\begin{document}
\maketitle

\noindent\emph{Source and licence.} The second variation of 2-spheres in the 2n-sphere. Version arXiv:2609.17887v1 (CC BY 4.0), distributed under the Creative Commons Attribution 4.0 International licence, as recorded in the arXiv licence metadata for this exact version and shown on the abstract page https://arxiv.org/abs/2609.17887v1. Full rights notice: licenses/arxiv-2609.17887-CC-BY-4.0.txt.\par\smallskip

\noindent\textit{Editorial scope.} Counted unit: the Prop whose source label is \texttt{prop:nullityinduction} in the article (source0.tex lines 1490--1501, source label \texttt{prop:nullityinduction}), delivered together with the article's own complete original proof of it (source0.tex lines 1503--1564, one \texttt{proof} environment of 62 lines). No mathematics is written, repaired or re-derived: statement and proof are the article's own text line for line. Required context retained uncounted: prop:commutident (lines 1177--1183); eq:qkJacobiidentities (lines 1440--1442); lem:ckdkcalc (lines 1450--1455). Every definition, display and label of those lines is retained unchanged. Omitted from this package and not redistributed: 1--1176, 1184--1439, 1443--1449, 1456--1489, 1502--1502, 1565--1828. Nothing inside a retained excerpt is omitted or altered: every retained body is delivered exactly as the article's lines are written, so no omission receipt of any kind accompanies this package. Citations: the retained excerpts carry no citation command at all, so no bibliography entry of the article is transcribed and no reference is redistributed. Reference adaptation: none was needed and none was made; every reference inside the delivered unit resolves inside it, so no unresolved reference or citation is produced. Printed slips kept literal: no printed word, symbol, hypothesis or number of the counted statement or of its proof was altered. Layout and class: the article's own class is \texttt{article} with packages including graphicx, inputenc, fontenc, amsmath, amsthm, amsfonts, amssymb, enumitem, ulem, natbib, url, slashed, bbm, tikz; the wrapper is the lane's pinned \texttt{amsart} editorial wrapper at 10pt on A4, which restates the theorem-like environments the retained text needs and adds the article's own macro definitions that the retained text needs (\texttt{\textbackslash del}, \texttt{\textbackslash delbar}, \texttt{\textbackslash llangle}, \texttt{\textbackslash rrangle}), with extra wrapper packages (bm, bbm, dsfont, xspace, cancel, mathtools, enumitem). The delivered numbering is the unit's own. Assets: no retained span contains a figure environment, a graphics-inclusion command, a PDF-page inclusion, a table or a TikZ picture; no image file is copied into this package, the package declares no asset dependency and no image file is redistributed with it. Licence: version arXiv:2609.17887v1 of this article is distributed under the Creative Commons Attribution 4.0 International licence, as recorded in the arXiv licence metadata for this exact version and shown on the abstract page https://arxiv.org/abs/2609.17887v1. A full rights notice is delivered at licenses/arxiv-2609.17887-CC-BY-4.0.txt. Characters: every retained excerpt is pure ASCII, so the delivered engine, XeLaTeX with its default Latin Modern text font, reports no missing character.
\begin{prop}\label{prop:commutident}
	For $0 \leq k \leq n-2,$ we have the identity
	\begin{equation*}
		(\overline{\partial_{\mathcal{E}_k}})^* \circ P_{k+1}^\dagger = -P_{k+1}^* \circ \partial_{\mathcal{E}_{k+1}},
	\end{equation*}
	where $P_{k+1}^*: \Omega^{1,0}(\mathcal{E}_{k+1}) \to \Gamma(\mathcal{E}_{k})$ is the formal $L^2$ adjoint of $P_{k+1}.$
\end{prop}

\begin{equation}\label{eq:qkJacobiidentities}
    \delbar^*_{\mathcal{E}_k} \delbar_{\mathcal{E}_k} W = P_{k+1}^* P_{k+1} W, \quad \del^*_{\mathcal{E}_k} \del_{\mathcal{E}_k} W = (P_k^\dagger)^* P^\dagger_k W.
\end{equation}

\begin{lemma}\label{lem:ckdkcalc}
    Let $f: \Sigma \to S^{2n}$ be a compact, linearly full, totally isotropic branched immersion. For $0 \leq k \leq n-2,$ we have
    \begin{equation*}
        c_k + d_k = (n-k-1)b_{k+1} + h^1(\mathcal{A}_k).
    \end{equation*}
\end{lemma}

\begin{prop}\label{prop:nullityinduction}
    Let $f: \Sigma \to S^{2n}$ be a compact, linearly full, totally isotropic branched immersion. For $0 \leq k \leq n-2,$
    \begin{enumerate}[label=(\alph*)]
        \item \begin{equation}
            \nu_k \geq \nu_{k+1} + \chi(\mathcal{A}_k) - (n-k-1)b_{k+1}
        \end{equation}
        \item
            \begin{equation}
                \nu_k \leq \nu_{k+1} + h^0(\mathcal{A}_k) + h^1(\mathcal{A}_k) + (n-k-1)b_{k+1}
            \end{equation}
    \end{enumerate}
\end{prop}

\begin{proof}
    Any pair $(W,Z) \in \Gamma(\mathcal{E}_k) \oplus \Gamma(\mathcal{E}_{k+1})$ satisfying the coupled system of equations
        \begin{equation}\label{eq:nullityextension}
            \delbar_{\mathcal{E}_{k}} W - P^\dagger_{k+1} Z = 0, \quad \del_{\mathcal{E}_{k+1}} Z + P_{k+1} W = 0.
        \end{equation}
    consists of a pair of $q_k$- and $q_{k+1}$-Jacobi fields. Indeed, by Proposition \ref{prop:commutident}, if $(W,Z)$ solves (\ref{eq:nullityextension}), then
        \begin{equation*}
            \begin{aligned}
                \mathcal{J}_k W &= \delbar_{\mathcal{E}_k}^* \delbar_{\mathcal{E}_{k}} W - P_{k+1}^* P_{k+1} W \\
                &= \left( \delbar_{\mathcal{E}_k}^* P^\dagger_{k+1} + P_{k+1}^* \del_{\mathcal{E}_{k+1}} \right)Z = 0.
            \end{aligned}
        \end{equation*}
        The calculation showing $\mathcal{J}_{k+1} Z = 0$ is similar.
    
        Let $U_{k} = \ker \mathcal{J}_k$ denote the space of $q_{k}$-Jacobi fields. For part (a), given $Z \in U_{k+1},$ we will try to construct $W \in \Gamma(\mathcal{E}_k)$ satisfying (\ref{eq:nullityextension}), while for part (b), given $W \in U_k,$ we will try to construct $Z \in \Gamma(\mathcal{E}_{k+1})$ satisfying (\ref{eq:nullityextension}).

        \begin{enumerate}[label=(\alph*)]
        \item We solve $\delbar_{\mathcal{E}_{k}} W - P^\dagger_{k+1} Z = 0$ first. In order to solve this equation, we need the class $[P_{k+1}^\dagger Z] \in H^1(\mathcal{E}_k)$ to vanish. Let $U'_{k+1}$ denote the subspace of $U_{k+1}$ for which this class vanishes. If $Z$ is a $q_{k+1}$-Jacobi field and if $\sigma \in \ker(\del_{\mathcal{E}_{k+1}}),$
        \begin{equation*}
            \begin{aligned}
                \llangle P_{k+1}^\dagger Z, P_{k+1}^\dagger \sigma \rrangle &= \llangle (P_{k+1}^\dagger)^* P_{k+1}^\dagger Z, \sigma \rrangle \\
                &= \llangle \del_{\mathcal{E}_{k+1}} Z, \del_{\mathcal{E}_{k+1}} \sigma \rrangle = 0,
            \end{aligned}
        \end{equation*}
        where we have used (\ref{eq:qkJacobiidentities}). Therefore, $P_{k+1}^\dagger U_{k+1}$ is orthogonal to $P_{k+1}^\dagger \ker(\del_{\mathcal{E}_{k+1}}),$ so the codimension of $U_{k+1}'$ in $U_{k+1}$ is at most $d_k.$

        For $Z \in U_{k+1}',$ choose the solution $W$ of $\delbar_{\mathcal{E}_{k}} W - P^\dagger_{k+1} Z = 0$ with $W \in \ker(\delbar_{\mathcal{E}_{k}} )^\perp$ and let $Y = \del_{\mathcal{E}_{k+1}} Z + P_{k+1} W.$ Proposition \ref{prop:commutident} and (\ref{eq:qkJacobiidentities}) give
        \begin{equation*}
            \begin{aligned}
                \del^*_{\mathcal{E}_{k+1}} Y &= \del^*_{\mathcal{E}_{k+1}} \del_{\mathcal{E}_{k+1}} Z + \del^*_{\mathcal{E}_{k+1}} P_{k+1} W \\
                &= (P_{k+1}^\dagger)^* P_{k+1}^\dagger Z + \del^*_{\mathcal{E}_{k+1}} P_{k+1} W \\
                &= \left( (P_{k+1}^\dagger)^* \delbar_{\mathcal{E}_{k}} + \del^*_{\mathcal{E}_{k+1}} P_{k+1} \right) W = 0,
            \end{aligned}
        \end{equation*}
        so $Y \in H^0(\mathcal{K} \otimes \mathcal{E}_{k+1}).$ Replacing $W$ by $W+\tau$ with $\tau \in H^0(\mathcal{E}_k)$ changes $Y$ by $P_{k+1} \tau.$ Let $U''_{k+1}$ denote the subspace of $U_{k+1}'$ for which the class of $Y$ in ${H^0(\mathcal{K} \otimes \mathcal{E}_{k+1})}/{P_{k+1} H^0(\mathcal{E}_k)}$ is $0.$ The codimension of $U''_{k+1}$ in $U'_{k+1}$ is at most $c_k.$ Given $Z \in U_{k+1}'',$ we can find $W' \in \Gamma(\mathcal{E}_k)$ satisfying both equations in (\ref{eq:nullityextension}), unique up to an element in $\ker(\delbar_{\mathcal{E}_k}) \cap \ker(P_{k+1}) = H^0(\mathcal{A}_k).$ Therefore, by Lemma \ref{lem:ckdkcalc},
        \begin{equation}
            \nu_k \geq \nu_{k+1} - c_k - d_k + h^0(\mathcal{A}_k) = \nu_{k+1} + \chi(\mathcal{A}_k) - (n-k-1)b_{k+1}.
        \end{equation}
        
        \item We solve $\del_{\mathcal{E}_{k+1}} Z + P_{k+1} W = 0$ first. In order to solve this equation, the orthogonal projection of $P_{k+1}W$ onto $H^0(\mathcal K \otimes \mathcal E_{k+1})$ must vanish. Let $V'_{k} < U_k$ denote the subspace of elements $W \in U_k$ for which this projection of $P_{k+1}W$ vanishes. If $W$ is a $q_k$-Jacobi field and if $\alpha \in H^0(\mathcal{E}_k),$ then
        \begin{equation*}
        \begin{aligned}
             \llangle P_{k+1} W, P_{k+1} \alpha \rrangle  &= \llangle P_{k+1}^* P_{k+1} W, \alpha \rrangle \\
             &= \llangle \delbar_{\mathcal{E}_k} W, \delbar_{\mathcal{E}_k} \alpha \rrangle = 0,
        \end{aligned}
        \end{equation*}
        where we have used (\ref{eq:qkJacobiidentities}). Therefore, $P_{k+1} U_k$ is orthogonal to $P_{k+1} H^0(\mathcal{E}_k),$ so the codimension of $V_k'$ in $U_k$ is at most $c_k.$

        For $W \in V_{k}',$ choose the solution $Z$ of $\del_{\mathcal{E}_{k+1}} Z + P_{k+1} W = 0$ with $Z \in \ker (\partial_{\mathcal E_{k+1}})^\perp$ and let $X = \delbar_{\mathcal{E}_{k}} W - P^\dagger_{k+1} Z.$ Proposition \ref{prop:commutident} and (\ref{eq:qkJacobiidentities}) give
        \begin{equation*}
            \begin{aligned}
                \delbar_{\mathcal{E}_k}^* X &= \delbar_{\mathcal{E}_k}^* \delbar_{\mathcal{E}_{k}} W - \delbar_{\mathcal{E}_k}^* P^\dagger_{k+1} Z \\
                &= P_{k+1}^* P_{k+1} W - \delbar_{\mathcal{E}_k}^* P^\dagger_{k+1} Z \\
                &= -\left( P_{k+1}^* \del_{\mathcal{E}_{k+1}} +  \delbar_{\mathcal{E}_k}^* P^\dagger_{k+1} \right) Z = 0,
            \end{aligned}
        \end{equation*}
        so $X \in H^1(\mathcal{E}_k).$ Replacing $Z$ by $Z + \gamma$ with $\gamma \in \ker (\partial_{\mathcal E_{k+1}})$ changes $X$ by $-P_{k+1}^\dagger \gamma.$ Let $V_{k}''$ denote the subspace of $V_k'$ for which the class of $X$ in ${H^1(\mathcal{E}_{k})}/({P^\dagger_{k+1} \ker(\del_{\mathcal{E}_{k+1}})})$ is $0.$ The codimension of $V_{k}''$ in $V_k'$ is at most $d_k.$ Given $W \in V_{k}'',$ we can find $Z' \in \Gamma(\mathcal{E}_{k+1})$ satisfying both equations in (\ref{eq:nullityextension}). The resulting map $V_{k}'' \to U_{k+1}$ has kernel contained in $H^0(\mathcal{A}_k),$ since if $Z' = 0$ then (\ref{eq:nullityextension}) gives $\delbar_{\mathcal{E}_k} W = 0$ and $P_{k+1} W = 0,$ so $W \in \ker(\delbar_{\mathcal{E}_k}) \cap \ker(P_{k+1}) = H^0(\mathcal{A}_k).$ Since $V_k''$ has codimension at most $c_k+d_k$ in $U_k,$ we have, using Lemma \ref{lem:ckdkcalc},
        \begin{equation*}
            \nu_{k+1} \geq \nu_k - c_k - d_k - h^0(\mathcal{A}_k) = \nu_k -h^0(\mathcal{A}_k) - h^1(\mathcal{A}_k) - (n-k-1)b_{k+1}.\qedhere
        \end{equation*}
    \end{enumerate}
\end{proof}

\end{document}
